\section{Introduction}
\label{sec:intro}

Retrieval-augmented generation (RAG) has converged on a deployment
template that few practitioners question: dense retrieval, then a
cross-encoder rerank, then per-passage compression or refinement,
then generation. Each stage is justified locally --- denser embeddings
pull the right passages, rerankers rescore them more carefully, and
refinement compresses each retained passage to its query-relevant
fragment. Each stage individually \emph{does} improve per-passage
relevance scores. The implicit and load-bearing assumption is that
because every stage helps locally, the composition helps globally:
better retrieval $\Rightarrow$ more faithful answers.

\paragraph{The phenomenon.} We report a counter-intuitive failure of
this assumption. Across $6{,}050+$ queries spanning SQuAD and
PubMedQA, two generator families (Mistral-$7$B and Llama-$3$-$8$B), and
four retrieval configurations, naive post-retrieval refinement
\emph{cuts} answer faithfulness from $0.7995$ to $0.6407$ on SQuAD and
introduces $10\%$ hallucination where there was none --- while leaving
per-passage similarity slightly higher
(Figure~\ref{fig:headline_frontier}). We name this the
\textbf{refinement paradox}.

\paragraph{The mechanism.} The culprit is not any individual passage
but a property of the retrieved \emph{set}. When refinement splinters
a broad passage into a narrow high-similarity fragment, it strips the
connective evidence the generator was relying on. The generator,
finding no longer-form context to ground its claims in, silently fills
the gap with parametric memory --- producing a fluent answer that no
longer follows from the supplied context. We call this property
\textbf{context coherence} and quantify it with the
\textbf{Context Coherence Score} (CCS): the mean minus the standard
deviation of pairwise cosine similarities between retrieved chunk
embeddings. CCS is generator-free, runs in $O(k^{2})$ in the embedding
dimension, and can be computed at retrieval time with no additional
forward pass.

\paragraph{The intervention.} If coherence is the active ingredient,
a retriever that \emph{preserves} coherence should recover faithfulness
without paying the latency cost of decoder-side reflection (e.g.,
Self-RAG~\citep{asai2024selfrag}). We instantiate this idea as
HCPC-v$2$, a coherence-gated refinement policy that declines to refine
when refinement would shatter the top-ranked block. HCPC-v$2$ recovers
SQuAD faithfulness to $0.7874$ with $0\%$ hallucination while
intervening on only $16.7\%$ of queries. Restraint, not algorithmic
sophistication, is the active ingredient.

\paragraph{Contributions.} We make four contributions.
\begin{enumerate}
  \item \textbf{The refinement paradox} (\S\ref{sec:paradox}): we
    document and characterize a failure mode in the standard RAG
    pipeline that holds across $5$ datasets, $5$ generator families
    ($7$B--$120$B), and $9$ ablation axes. To our knowledge this is
    the first systematic demonstration that better per-passage
    retrieval can make answer faithfulness \emph{worse} on standard
    short-answer QA.
  \item \textbf{Context Coherence Score (CCS)} (\S\ref{sec:method}): a
    generator-free retrieval-time diagnostic that predicts the
    paradox. CCS separates high-coherence sets ($0\%$ hallucination)
    from low-coherence sets that precede generation failures; the
    model's own self-reported confidence is significantly correlated
    with CCS (Pearson $r{=}0.36$, Spearman $\rho{=}0.48$, both $p{<}0.01$),
    cross-validating that the LLM internally tracks coherence.
  \item \textbf{HCPC-v$2$ as a deployment policy} (\S\ref{sec:method}):
    a coherence-gated retriever that recovers faithfulness with $0\%$
    hallucination while intervening on $1$ in $6$ queries. Adds one
    pairwise-similarity computation to the request path; no
    fine-tuning, no labeled data, no decoder reflection.
  \item \textbf{Robustness across $7$ axes} (\S\ref{sec:robustness}):
    multi-seed variance, hierarchical retrieval (RAPTOR), prompt
    templates, sub-chunk size, frontier scale ($70$B/$120$B), reranker
    family, quantization, and sampling temperature. The paradox
    persists where the coherence account predicts and is honestly
    qualified where it does not.
\end{enumerate}

\paragraph{Released artifacts.} To make the result fully usable: the
benchmark is on Zenodo with permanent DOI
\href{https://doi.org/10.5281/zenodo.19757291}{\texttt{10.5281/zenodo.19757291}};
a HuggingFace Dataset
(\href{https://huggingface.co/datasets/saketmgnt/context-coherence-bench}{\texttt{saketmgnt/context-coherence-bench}})
provides \texttt{load\_dataset}-able access to $5$ configurations
including the adversarial cases; an interactive Gradio demo lives at
\href{https://huggingface.co/spaces/saketmgnt/sakkk}{huggingface.co/spaces/saketmgnt/sakkk};
and the standalone CCS metric ships as a $\texttt{pip install}$-able
package and as a drop-in LangChain
\texttt{CoherenceGatedRetriever}.

\paragraph{Roadmap.} \S\ref{sec:bg} positions the work against prior
RAG literature. \S\ref{sec:setup} describes the experimental design.
\S\ref{sec:paradox} presents the paradox itself.
\S\ref{sec:theory} gives an information-theoretic account.
\S\ref{sec:method} introduces CCS and HCPC-v$2$.
\S\ref{sec:robustness} bundles seven robustness checks.
\S\ref{sec:discussion} discusses limits, including three regimes
where CCS is \emph{not} predictive. \S\ref{sec:conclusion}
concludes.
